\chapter{Conclusion}

\section{Résumé du Stage}
    
    Dans ce rapport, les principales activités et réalisations effectuées durant mon stage ont été présentées. Ces travaux visaient à améliorer le processus MoCapIA, allant de la calibration jusqu'à la reconstruction 3D d'un modèle humain en passant par l'estimation 2D sans marqueurs et la triangulation des données pour des applications cliniques, sportive ou industrielle. Quatre anciens stagiaires avaient auparavant travaillé sur le projet. Ils ont développé la majorité des fonctions de base du processus et testé les différentes technologies à dispositions (ex : modèle d'estimateur de pose 2D) tout en respectant les conventions de programmations.
    
    Il a donc été facile durant la première phase, de comprendre le code existant et se familiariser avec l'Interface Utilisateur. Avec l'aide du précédent stagiaire et de mon tuteur de stage, une démo a été réalisée en début de stage ce qui m'a permis d'observer les différentes étapes et surtout d'identifier les points faibles du projet MoCapIA. Par la suite, mon tuteur de stage m'a incité à conduire une veille scientifique en se reposant sur une recherche bibliographique approfondie dans les domaines de la capture de mouvement et de l'estimation de pose afin d'identifier les avancées récentes. J'ai rapidement compris que cette étape était essentielle pour comprendre les enjeux, mais notamment pour identifier les éventuelles améliorations.
    
    Dans un deuxième temps, l'objectif était d'identifier les points faibles du processus sur la base de mesures expérimentales. Pour ça, on a mis en place une série d'expériences :
    \begin{itemize}
        \item \textbf{TEST 0} afin de comprendre l'impact du damier sur la qualité de la calibration.
        \item \textbf{TEST 1} sur l'homogénéisation de la calibration en mettant de côtés les erreurs d'estimation 2D.
        \item \textbf{TEST 2} a permis de comparer RTMPose, l'estimateur de pose 2D utilisé dans MoCapIA, avec une méthode de référence avec marqueur (Vicon).
    \end{itemize}
    Ces différents tests nous ont permis de conclure que l'estimation 2D des points caractéristiques était le maillon faible. 
    
    Dans un troisième temps, une importante étape de recherche bibliographique a été réalisée afin de remédier aux problèmes mis en lumière ci-dessus. J'ai donc implémenté de nouvelles méthodes pour affiner et corriger les points 2D détectés. En s'inspirant d'un projet similaire (Pose2Sim) qu'on a d'abord comparé avec MoCapIA, 3 étapes ont été retenues et adaptées à notre processus. Le filtrage a permis de supprimer les valeurs incohérentes et de supprimer le bruit lié aux erreurs d'estimation 2D. Les deux autres étapes ont permis, en plus de figer la distance des membres en appliquant un squelette au modèle 3D,d'augmenter le nombre de marqueurs afin de permettre une analyse biomécanique bien plus complète. Chaque étape a été testée et validée individuellement avant d'être intégrée dans le processus global. Des tests finaux ont été réalisés pour valider l'ensemble des améliorations apportées. Les résultats ont montré une nette amélioration de la qualité de la reconstruction 3D, confirmant l'efficacité des nouvelles étapes implémentées. 
    
    Enfin, grâce à une veille technologique, nous avons pu tester un nouvel outil de capture 3D de mouvement basée uniquement sur une caméra et l'IA SAM3D de Meta. En effet, même si son application sur une vidéo reste pour l'instant complexe, SAM3D a ouvert une nouvelle voie dans le domaine de la capture de mouvement avec des résultats très prometteurs. Il est possible de modéliser le mouvement d'un être humain en 3D avec seulement une caméra et sans avoir besoin de calibration, on se rapproche encore plus des objectifs de moindres coûts et de simplicité. Entre temps, j'ai ajouté des fonctionnalités secondaires au logiciel ; pour permettre de gérer quelques paramètres comme la possibilité de choisir précisément quelles parties lancées (ex : pouvoir choisir quel filtre appliqué avec quels paramètres) qui n'a pas pu être ajouté dans l'interface utilisateur. Les étapes principales ajoutées, elles, ont par ailleurs été intégrées dans l'interface grâce à la deuxième stagiaire Yiyang.

\section{Analyse réflexive}
    
    \paragraph{Analyse Critique}

        Cette partie va permettre d'effectuer une conclusion critique et personnelle sur le travail accompli et ces six mois passé au BMBI, en soulignant ce qui a fonctionné et ce qui ne l'a pas été. Je vais aussi explorer les points positifs et négatifs de cette mission. Enfin, je vais mettre en lumière les compétences que j'ai pu développer durant ce stage et leur cohérence ou non avec mon projet professionnel.

        Premièrement le sujet du stage m'a beaucoup plu. En effet, la capture de mouvement est un domaine qui m'intéresse depuis longtemps, j'avais d'ailleurs effectué un projet de fin de cursus en classe préparatoire sur le thème de la capture de trajectoire de balle. Ce projet reprenait en grande partie les notions de calibration et de triangulation, mais aussi de traitement d'images et de détection par une IA qu'on a utilisé dans le projet MoCapIA. J'ai donc pu approfondir mes connaissances dans ces domaines et découvrir de nouvelles notions comme la dimension biomécanique et les technologies d'estimations de pose humaine.

        Deuxièmement, la dimension "recherche" du stage a été vraiment une bonne expérience. Cela m'a permis de me familiariser avec la recherche bibliographique scientifique en anglais et de comprendre l'importance de cette étape dans le développement d'un projet innovant. Je pense avoir acquis une bonne méthodologie de recherche documentaire qui me sera utile dans le futur. Que ce soit dans le domaine de la recherche ou de l'entreprise, la rigueur et la nécessité de valider l'ensemble de ses choix par des sources fiables est, à mon sens, un point essentiel.

        Troisièmement, ce stage a été effectué à deux. En effet, avec Yiyang Huang, une étudiante GI04, elle aussi en TN09, nous avons partagé le même tuteur de stage et avons travaillé sur le même projet MoCapIA. Cela a impliqué la mise en place d'une coordination et d'une collaboration efficace pour la réalisation des tâches, ce qui au final a été l'un des points positifs du stage. En effet, ce travail collaboratif pour la conception d'un logiciel et les tests menés avec le groupe de TX a permis d'introduire à plus grande échelle la notion de travail en équipe sur quelque chose de concret. Nous avons utilisé des outils de gestion de version (Git / GitLab pour le stockage) et de communication (réunions hebdomadaires). Le fait d'avoir une autre personne avec qui échanger sur les problématiques rencontrées est aussi un point vraiment important. 
        
        Les réalisations qui ont été exposés dans ce rapport ont permis d'atteindre les objectifs fixés au début du stage. Les étapes de test ont été concluantes et m'ont servis tout au long du stage pour orienter mes recherches. On peut toutefois critiquer certains aspects des expériences. En effet, lors du \textbf{TEST 1 - Analyse de l'homogénéisation de la calibration}, on triangule quand même les résultats pour calculer l'erreur de calibration. Lors des clics effectués pendant l'expérience, on obtient toujours une erreur de quelques pixels ce qui implique une légère approximation quand on applique l'algorithme de triangulation. Une étude de sensibilité pourrait compléter l'analyse. En affirmant que l'on mesure uniquement l'erreur de calibration, on néglige cette erreur de triangulation et on émet l'hypothèse qu'elle reste négligeable par rapport à l'erreur de calibration. Cependant, cette erreur est très faible comparée à l'erreur de calibration, on considère donc cette hypothèse comme acceptable. Je rajouterai que lors du test entre le système optoélectronique Vicon et le système MoCapIA, on considère le Vicon comme la référence absolue. Il est pourtant important de noter que ce système souffre aussi d'erreurs, notamment en provenance des approximations humaines lors du placement de marqueurs. Les erreurs observées lors de la conclusion sont donc à nuancer \cite{MarkerlessforBiomecha}. On peut aussi nuancer la conclusion de la comparaison entre MoCapIA et SAM3D. En effet, la comparaison de SAM a été faite avec la version de MoCapIA sans les corrections biomécaniques.

        Des difficultés ont été rencontrées durant le stage. La plus importante a été la recherche bibliographique. En effet, c'était quelque chose d'assez nouveau. Au début, j'étais incapable de cibler les bonnes informations dans les articles et j'utilisais des assistants conversationnels pour m'aider à résumer les articles. Cette approche n'était pas la bonne, mais au fur et à mesure, j'ai pu m'améliorer et apprendre à extraire les informations importantes par moi-même. Mon tuteur de stage m'a aussi beaucoup aidé en me guidant dans mes recherches et en me donnant des pistes à explorer tout au long du stage. Une autre difficulté rencontrée a été la compréhension du code. En effet, étant un projet de stage depuis deux ans, le logiciel contenait plusieurs dizaines de fichiers et plusieurs milliers de lignes de code. Même si les conventions de programmations étaient respectées, il a fallu un certain temps pour se familiariser avec l'architecture du code et le dialogue entre les différentes fonctions et l'interface utilisateur.

        Finalement ce stage m'a permis de développer plusieurs compétences et d'affiner mon projet professionnel. En matière d'intelligence artificielle, j'ai pu approfondir mes connaissances sur les modèles de deep learning, avec l'utilisation des bibliothèques PyTorch et Tensorflow.
        J'ai pu renforcer mes compétences en traitement d'images et en vision par ordinateur, notamment à travers les étapes de calibration, d'estimation de pose 2D et de la reconstruction 3D. J'ai pu m'initier à la programmation orientée objet en Python (habitué au C++) en travaillant sur ce projet 100\% codé en python. Le concept est le même, mais j'ai pu remarquer plusieurs différences, car Python est plus souple que le C++. Pour ce qui est des algorithmes de traitement de résultats, surtout pour des tests statistiques et graphiques, j'ai majoritairement utilisé MATLAB, ce qui me permet aujourd'hui d'être encore plus à l'aise avec ce langage. L'aspect mathématique du projet m'a beaucoup plu, particulièrement avec les notions de géométrie projective et d'algèbre linéaire. L'équipe BioMov-E accueille plusieurs doctorants et chercheurs internationaux donc chaque réunion hebdomadaire se déroulait en anglais ; cela m'a permis d'améliorer mon niveau d'anglais scientifique. 
        Je n'exclus pas la possibilité de poursuivre dans la recherche, même si je compte orienter ma recherche de stage TN10 en entreprise. Je pense m'orienter vers la filière INES pour la suite de mon cursus, car j'ai apprécié la vision par ordinateur et j'aimerais l'appliquer à la robotique ou à des systèmes autonomes.


    \paragraph{Impact sociétal / éthique}

        Ce projet s'inscrit dans un contexte dans lequel la technologie se retrouve au service de la santé. Il a pour but d'améliorer les diagnostics des professionnels de santé en leur fournissant des outils plus précis sans pour autant augmenter considérablement les coûts. En rendant la capture de mouvement plus accessible, on permet à un plus grand nombre de patients de bénéficier d'analyses détaillées de leurs mouvements. Utilisant des modèles d'intelligence d'artificielle, l'utilisation de cette technologie pourrait inquiéter le patient sur les données recueillies et utilisées. Pourtant, ces données se réduisent pour l'instant à des informations de base comme la taille, le poids, l'âge et le sexe ; cela reste donc relativement peu intrusif. De plus, le logiciel est aussi destiné aux sportifs de hauts niveaux ou aux entraîneurs, leur permettant d'analyser les performances et de prévenir les blessures grâce à une analyse biomécanique détaillée.

        Le projet a donc plusieurs intérêts sociétaux. Les impacts environnementaux sont, eux aussi, limités. En effet, en utilisant des caméras classiques et en réduisant le nombre de capteurs nécessaires, on diminue la consommation de ressources matérielles. Même si beaucoup d'énergie ont été nécessaires pour entraîner les modèles d'IA que l'on utilise, l'impact sur la durée totale de vie reste vraiment faible.

    \paragraph{Perspectives}

        L'écriture de ce rapport a permis de faire le point sur les différentes étapes du projet et d'identifier des pistes d'améliorations futures. Plusieurs axes de développement peuvent être envisagés pour continuer à améliorer le processus MoCapIA. En suivant de près les avancées dans le domaine de l'IA, on pourrait intégrer de nouveaux modèles d'estimation de pose 2D encore plus performants. Même si RTMPose reste le modèle le plus adapté à la fin de ce stage, de nouvelles technologies font leur apparition. C'est le cas pour le modèle SAM3D qui a montré des résultats prometteurs. Si la prochaine équipe réussie à comprendre plus en détails cette technologie et réussi à l'adapter à des vidéos complètes, on pourrait répondre encore mieux aux objectifs initiaux du projet MoCapIA.
        L'interface utilisateur pourrait aussi être amélioré en intégrant plus de paramètres de contrôle pour les différentes étapes du processus et en s'inspirant des logiciels de santé et de gestion de patient.
        Une façon plus simple de calibrer avec les caméras pourrait être implémentée. Avec une disposition de caméras aux quatre angles de la zone de capture, la calibration avec un damier devient compliquée. L'utilisation d'objets de calibration en trois dimensions pourraient simplifier cette étape \cite{OtherCalibMethods}.